\chapter{Introduction}
\label{ch:introduction}

A transpiler between $k$ languages needs a rule for each ordered pair of them, and so $k(k-1)$ rules. A
transpiler through one intermediate form needs $2k$ maps: one into the form and one out of it for each language.
The orior transpiler takes the second design. Its intermediate form is gsm, the relational assembly of gnascor
\cite{src:gnascor}. A gsm statement does not name an operation. It states the cases of a relation, as
\texttt{1,1->2}, and a machine either answers those cases or does not.

The claim of this research paper is narrow and exact. Each map into gsm and out of it is a finite table, and each
entry of each table carries a proof. The proofs are Lean 4 theorems \cite{src:lean4} over the tables as the
repository stores them. Nothing in the proofs holds a second copy of a table. A tool writes the tables into Lean
data, and the kernel of Lean checks each theorem against that data.

\section{The transform}

Let $\Sigma_L$ be the symbols of a language $L$, let $B$ be the bridge words of the transpiler, and let $A$ be the
anchors of the ladder. Three finite partial maps make the transform:
\begin{align*}
  \tau_L &: \Sigma_L \rightharpoonup B, &
  \alpha &: B \rightharpoonup A, &
  R &: A \to \bigl(W_n \times W_n \rightharpoonup W_n\bigr).
\end{align*}
Here $W_n$ is the set of words of $n$ places, $\tau_L$ is the target table of $L$, $\alpha$ is the anchor column
of the transpiler table, and $R_a$ is the relation of the anchor $a$.

\begin{definition}
A chart is a target table $\tau_L$. The atlas is the set of charts $\{\tau_L\}$ together with $\alpha$ and $R$.
Every chart maps into the one set $B$, and $\alpha$ and $R$ carry $B$ to the one set of relations.
\end{definition}

The meaning that the atlas gives a symbol $\sigma$ is the relation $R_{\alpha(\tau_L(\sigma))}$. The gsm of
$\sigma$ is the list of cases of that relation in ladder.h. A new language adds one chart, and every chart it meets
shares the same relations.

\begin{definition}
The multi-cardinality of the atlas is the path
\[
  \text{any} \;\to\; \text{atlas} \;\to\; \text{gsm} \;\to\; \text{atlas} \;\to\; \text{any}.
\]
A symbol of any language enters the atlas through its chart, projects to gsm, and leaves the atlas through the
chart of any language.
\end{definition}

The path through gsm is the matrix projection of the title. The projection has four cardinalities.
\begin{itemize}
  \item 1:1. One symbol of one language projects to one symbol of another, as C's $+$ to Python's masked $+$.
  \item 1:many. One symbol projects to a symbol of each of many languages, each through its own chart.
  \item many:many. Many symbols reach one relation and leave as many symbols. C's $+$ and $+{=}$ reach add, and
        add leaves as Python's $+$ and $+{=}$.
  \item n:many. The charts of $n$ languages project through the one gsm to the charts of many languages.
\end{itemize}
One term can be in more than one category. The atlas keeps the intent of a term until it transcribes the term's
meaning to gsm. Each lift is a projection in and a projection out, whatever its cardinality.

A lift of $\sigma \in \Sigma_L$ into a second language $M$ is a symbol $\sigma' \in \Sigma_M$ with
$\tau_M(\sigma') = \tau_L(\sigma)$. Chapter~\ref{ch:targets} shows that a lift keeps the meaning of $\sigma$
wherever both sides carry a proof.

\section{What the proofs cover}

The proofs fall into four groups, one to a chapter.
\begin{itemize}
  \item Chapter~\ref{ch:ladder}: the ladder relations, their cases, and their readings at every width and sign.
  \item Chapter~\ref{ch:gsm}: the state-transition matrices of gsm, and the transpiler form of each mnemonic.
  \item Chapter~\ref{ch:g}: the high order face \texttt{.g}, and the compile of each operator to gsm.
  \item Chapter~\ref{ch:targets}: the target tables of C and Python, and the lift between any two languages.
\end{itemize}
Chapter~\ref{ch:integration} gives the procedure that checks the whole transform at once, and the limits that
remain. Every theorem below has the name that it has in the Lean sources under \texttt{theory/lng}.
